\documentstyle[%


righttag%


,11pt%


,amssymb%


,russian%


,izvru%


]{amsart}





%       Definitions





\newcommand{\eps}{\varepsilon}


\newcommand{\de}{\delta}


\newcommand{\ovl}{\overline}


% \newcommand{\ov}{\overline}


\newcommand{\si}{\sigma}


\newcommand{\om} {\omega}


\newcommand{\ovo}{\ov{\omega}}


\newcommand{\vf}{\varphi}


% \newcommand{\wt}{\widetilde}


% \newcommand{\wh}{\widehat}


\newcommand{\no}{\nonumber}


\newcommand{\be}{\begin{equation}}


\newcommand{\ee}{\end{equation}}


\renewcommand{\baselinestretch}{1.05}


\theoremstyle{plain} %% This is the default


\begingroup


\theorembodyfont{\it}


\newtheorem{thmm}{\hskip\parindent Теорема}[section]


\endgroup





\numberwithin{equation}{section}





\newcommand{\pa}{\partial}


\newcommand{\dxx}{\frac{\pa^k}{\pa x^k}}


\newcommand{\dxt}{\frac{\pa^{k+k_0}}{\pa x^k \pa t^{k_0}}}








\setcounter{page}{98}


\god{1997}


\nomer{4~(419)}


\udk{519.633}


\author{\it Г.И. ШИШКИН, И.В. ЦЕЛИЩЕВА}


\title{МЕТОД ДЕКОМПОЗИЦИИ ДЛЯ СИНГУЛЯРНО ВОЗМУЩЕННЫХ\\


КРАЕВЫХ ЗАДАЧ С ЛОКАЛЬНЫМ ВОЗМУЩЕНИЕМ НАЧАЛЬНЫХ \\


УСЛОВИЙ. УРАВНЕНИЯ С КОНВЕКТИВНЫМИ ЧЛЕНАМИ}





\thanks{Работа выполнена при финансовой поддержке Российского фонда


фундаментальных исследований (код проекта 95-01-00039),


а также частично при финансовой поддержке Нидерландской исследовательской 


организации NWO (грант \No 047.003.017).}


\date{}


\pagestyle{plain}


\begin{document}





\maketitle





Рассматривается задача Дирихле для сингулярно возмущенных параболических 


уравнений в случае одной пространственной переменной (на отрезке). 


Начальное условие задачи имеет единичное локальное возмущение 


конечной амплитуды на узкой подобласти вблизи $x=0$ ширины $2\delta$. 


Возмущающий параметр $\eps^2$ --- коэффициент при старших


производных уравнения, а также параметр $\delta$ могут принимать 


произвольные значения из полуинтервалов $(0,1]$ и $(0,d]$ соответственно,


где $2d$ --- длина отрезка. 


При $\eps=0$ параболическое уравнение


вырождается в гиперболическое уравнение первого порядка, содержащее 


производные по пространственной и временной переменным. 


Такого типа


задачи возникают при моделировании процессов распространения тепла 


(при $t>t_0>0,$ $t_0=\eps^{-2}\delta^2$) в случае сосредоточенных 


мгновенных источников.





Показано, что для таких задач в случае классических разностных аппроксимаций


не существует прямоугольных кусочно равномерных сеток, на которых решение


разностной схемы сходилось бы равномерно относительно параметров $\eps$


и $\delta$ (или, короче, $(\eps,\delta)$-равномерно). 


Для указанных краевых задач с использованием метода аддитивного


выделения особенностей и подвижных сгущающихся 


(в окрестности переходного слоя) сеток, 


узлы которых расположены вдоль характеристики предельного 


уравнения, строятся монотонные разностные схемы, сходящиеся 


$(\eps,\delta)$-равномерно. 





\section*{Введение}





Решения краевых задач для сингулярно возмущенных уравнений


в случае гладких начальных условий, изменяющихся


на конечную величину в узкой области, обладают ограниченной гладкостью.


Производные решения неограниченно возрастают, когда величина $\delta$ ---


полуширина области резкого изменения начальных данных --- и (или)


возмущающий параметр $\eps$ стремятся к нулю. 


Наличие локальных возмущений приводит к появлению внутренних


(переходных) слоев при малых значениях параметров $\eps$ и $\delta$.


Заметим, что задача является сингулярной даже при $\eps=1$; при стремлении


параметра $\delta$ к нулю гладкость решения ухудшается. Эта ограниченная


гладкость решения краевой задачи вызывает затруднения при ее численном 


решении (см., напр., \cite{Ba}--\cite{Sh1}). При малых значениях параметров 


$\eps$ и $\delta$ ошибки приближенных решений, получаемых с использованием


классических разностных схем, становятся соизмеримыми с искомым решением.


В связи с этим ставится задача разработки специальных разностных


схем, ошибка решений которых не зависит от величины параметров


$\eps$ и $\delta$, т.е. разностных схем, сходящихся $(\eps,\delta)$-равномерно.





Для численного решения краевой задачи при малых значениях 


параметров $\eps$ и $\delta$


применяется метод аддитивного выделения особенностей (его описание см., 


напр., в (\cite{Mar}, гл.\,4, с.\,208)):


отдельно вычисляется главный член сингулярной части


решения, порождаемой локальным возмущением. При построении разностных


схем используются классические аппроксимации краевой задачи. 


Для вычисления решения при не слишком малых значениях параметров, а также


гладкой части решения применяются равномерные прямоугольные сетки;


при вычислении сингулярной части решения используются как прямоугольные 


сетки, так и подвижные ``характеристические'' сетки, узлы которых расположены


вдоль характеристики предельного уравнения, проходящей через точку $(0,0).$


На таких сетках решение разностной схемы сходится 


$(\eps,\delta)$-равномерно. На базе этого подхода может быть 


применен метод декомпозиции области в окрестности локального


возмущения начального условия (описание метода см., напр., в \cite{Sh2}, где


рассматривались уравнения без конвективных членов). Отметим, что 


метод аддитивного выделения особенностей в случае сингулярно возмущенных


краевых задач для уравнений с конвективными членами не применялся.








\section{Постановка задачи}





{\bf 1.1.} На отрезке $D=\{x:-d<x<d\}$ рассмотрим краевую задачу


для параболического уравнения 


\begin{gather} \label{1.1}


Lu \equiv \eps^2a(x,t)\frac{\pa^2 u}{\pa x^2} +


b(x,t)\frac{\pa u}{\pa x}-c(x,t)u-p(x,t)\frac{\pa u}{\pa t}=


f(x,t),\quad (x,t)\in{G}, \\ 


u(x,t)=\vf(x,t),\quad (x,t)\in S.\notag


\end{gather}


Здесь 


$G=D\times(0,T],$ $S=S(G)=\ovl{G} \setminus G$,


коэффициенты уравнения и правая часть, а также функция $\vf(x,t)$ 


являются достаточно гладкими и ограниченными на множестве $\ovl{G}$ и на


сторонах области $G$ соответственно, функция $\vf(x,t)$ непрерывна на $S$,


более того, 


$$


0<a_0\le a(x,t)\le a^0,\quad p(x,t)\ge p_0>0,\quad 


c(x,t)\ge 0,\quad (x,t)\in G;


$$


параметр $\eps$ принимает произвольные значения из полуинтервала $(0,1].$





Функция $\vf(x,t)$ зависит от параметра $\delta$, 


принимающего произвольные значения из полуинтервала $(0,d]$. 


При $t=0$ функция $\vf(x,t)=\vf(x,t;\delta)$, являясь гладкой и ограниченной


на $\ovl{D}$, изменяется на конечную величину в $\delta$-окрестности 


множества $\Gamma=\{x=0\}$. Пусть


$\vf(x,t)=\vf_0(x),$ $(x,t)\in S_0$, где 


$S_0=\{(x,t):x\in\ovl{D},\;t=0\}$ --- нижнее основание множества $\ovl{G}$, 


$S=S_0\cup S_1$.


Для функции $\vf(x,t)$ справедливы оценки 


\begin{align}\label{1.2}


&\left|\frac{\pa^{k_0}}{\pa t^{k_0}}\varphi(x,t)\right|\le M, \quad


(x,t)\in S_1,\\


%


&\left|\frac{\pa^k}{\pa x^k}\varphi(x,t)\right|\le M,\quad 


|x|\ge\delta,\notag\\


%


&\left|\frac{\pa^k}{\pa x^k}\varphi(x,t)\right|\le M\delta^{-k},\quad


|x|<\delta,\quad (x,t)\in S_0.\notag


\end{align}


Здесь и ниже через $M$ (или $m$) обозначаем достаточно большие (малые)


положительные постоянные, не зависящие от величины параметров $\eps$ и $\delta$.


В случае сеточных задач эти постоянные не зависят и от параметров шаблонов


используемых разностных схем.





Требуется найти решение задачи (1.1).


Под решением задачи понимается функция $u\in C^{2,1}(G)\cap C(\ovl{G})$, 


ограниченная на $\ovl{G}$, удовлетворяющая дифференциальному уравнению на $G$ и граничному


условию на $S$. Предполагаем выполнение условий согласования, обеспечивающих


гладкость решения при каждом наборе значений параметров $\eps$ и $\delta$.





При стремлении параметра $\eps$ к нулю в окрестности множества $S_1$


появляются пограничные слои, а при $\delta<\eps$ --- переходный слой


в окрестности множества $S^*=\{x=0\}\times(0,T]$. Заметим, что задача 


является сингулярной и при $\eps=1$; при стремлении параметра $\delta$ к нулю 


гладкость решения ухудшается: при фиксированном значении $\delta$ 


решение задачи достаточно гладко, однако его производные


неограниченно возрастают при $\delta\to 0$.





Такого типа задачи возникают при моделировании конвективных 


процессов теплопередачи (или диффузии вещества) в теле, 


составленном из частей, имеющих разную начальную температуру 


(концентрацию вещества). Для таких процессов


рассматриваются моменты времени, большие некоторого момента


$t_0=\eps^{-2}\delta^2$; выбором параметра $\delta=\delta(\eps)$


величина $t_0$ может быть сделана сколь угодно малой. К исследованию


задачи (1.1) с граничными условиями, удовлетворяющими (1.2), 


приводят, например, задачи о распространении тепла (при $t\ge t_0>0$) в


случае сосредоточенных мгновенных источников


(\cite{Ti}, гл.\,III, \S\,2, с.\,202).





{\bf 1.2.} Отметим проблемы, возникающие при численном решении задачи 


классическими методами. На отрезке $D=\{x:-1<x<1\}$


рассмотрим модельную краевую задачу для регулярного уравнения 


теплопроводности ($\eps=1$)


\begin{gather}\label{1.3}


Lu\equiv \left\{ \frac{\partial^2}{\partial x^2}+


\frac{\partial}{\partial x}-


\frac{\partial}{\partial t} \right\} u(x,t)=0,\quad (x,t)\in G,


\tag{1.3а}\\


%


u(x,t)=0,\quad (x,t)\in S_1,\quad u(x,t)=\vf_0(x),\quad (x,t)\in S_0.


\tag{1.3б}\\


%


\intertext{Функция $\vf_0(x)=\vf_0(x;\delta)$ удовлетворяет условиям}


%


\vf_0(x)=0,\quad |x| \ge \delta,\quad


\left| \frac{d^k}{dx^k} \vf_0(x) \right| \le M\delta^{-k},\quad


x\in\ovl{D}.\tag{1.3в}


\end{gather}





Для решения этой задачи используем классическую разностную схему 


\begin{gather}\label{1.4}


\tag{1.4}


\begin{split}


&\Lambda z\equiv \{\delta_{x\,\ovl{x}}+\delta_x-


\delta_{\ovl{t}}\}z(x,t)=0,\quad (x,t)\in G_h,\\


%


&z(x,t)=0,\quad (x,t)\in S_{1h},\quad z(x,t)=\vf_0(x),\quad (x,t)\in S_{0h}.


\end{split}


\end{gather}


Здесь $\ovl{G}_h$ --- прямоугольная равномерная сетка с числом узлов 


$N+1$ и $N_0+1$ по $x$ и $t$ соответственно; 


$\delta_{x\,\ovl{x}}\,z(x,t)$ и


$\delta_{x\,}z(x,t)$, $\delta_{\ovl{t}}\,z(x,t)$ --- вторая и первые 


(вперед и назад) разностные производные.





Перейдем к новой системе координат, в которой производные начальной функции 


ограничены равномерно относительно параметра $\delta$. 


Полагаем $\xi=\delta^{-1}x$, $\tau=\delta^{-2}t$.


Функция $\wt{\vf}_0(\xi)=\vf_0(x(\xi))$ удовлетворяет условиям 


$$


\wt{\vf}_0(\xi)=0,\quad |\xi| \ge 1,\quad \left| \frac{d^k}{d\xi^k}


\wt{\vf}_0(\xi) \right| \le M,\quad \xi\in\wt{D}.


$$


Отметим, что в новых переменных $\xi,$ $\tau$ шаги сеток, порождающих сетку


$\ovl{\wt{G}}_h$, будут определяться соотношениями 


$h_{\xi}=2\delta^{-1}N^{-1}$, $h_{\tau}=\delta^{-2}TN_0^{-1}$.





Выполнив анализ задачи в новых независимых переменных $\xi,$ $\tau$ 


(см, напр., \cite{Sh3}), устанавливаем, что для любого сколь угодно большого 


числа $N$ найдется такое значение $\delta=\delta(N)=N^{-1}$, 


при котором для 


погрешности решения будет выполняться неравенство


$$


\max\limits_{\ovl{\wt{G}}_h}\, |


\wt{u}(\xi,\tau)-\wt{z}(\xi,\tau) | \ge m>0


$$


при любых значениях $N_0$. Возвращаясь к исходным переменным,


получаем


$$


\max\limits_{\ovl{G}_h}\, |u(x,t)-z(x,t) | \ge m>0


$$


для любых $N,$ $N_0$ при $\delta=\delta(N)$.


Таким образом, приходим к следующему утверждению.





\begin{thmm} \label{th0}


Функция $z(x,t),$ $(x,t)\in\ovl{G}_h$, --- решение классической разностной схемы


{\em (1.4)} на равномерных сетках --- не


сходится к решению краевой задачи 


{\em (\ref{1.3})} $\delta$-равномерно при $N, N_0\to \infty$.


\end{thmm}





Так как основной является операторная формулировка проблемы, то и в 


случае более общей задачи с особенностями вида (1.1) мы приходим к


проблеме разработки специальных разностных схем, погрешность решения


которых не зависела бы от значений возмущающих параметров.





Для простоты предполагаем, что данные краевой задачи 


удовлетворяют условиям, при которых пограничные слои не возникают. 


Построение схем при наличии пограничных слоев рассматривается, например,


в \cite{Sh1}, \cite{Sh4}.


 


\section {Априорные оценки. Вспомогательные построения}





Для решения краевой задачи (1.1) и его производных получим ряд оценок 


на основе асимптотического представления решения, предполагая выполненными


условия (1.2). При выводе оценок применяется техника из


\cite{Sh1}, \cite{Sh4}, \cite{Sh5} 


с использованием мажорантных функций.





Представим решение задачи (1.1) в виде суммы функций


\begin{equation}\label{2.1}


u(x,t)=U(x,t)+V(x,t),\quad (x,t)\in\ovl{G}.


\end{equation}


Функции $U(x,t)$ и $V(x,t)$ --- гладкая (регулярная) и сингулярная части 


решения --- являются решениями задач 


\begin{alignat}{2}\label{2.2} 


L_{(\ref{1.1})}U(x,t) & = f(x,t),\qquad &(x,t)\in G,\notag\\


%


U(x,t) & = \vf^1(x,t),\qquad &(x,t)\in S;\tag{2.2а}\\


%


L_{(\ref{1.1})}V(x,t) & = 0,\qquad &(x,t)\in G,\notag\\


V(x,t) &= \vf^2(x,t),\qquad &(x,t)\in S.\tag{2.2б}


\end{alignat}


Здесь $\vf^i(x,t)$, $i=1,2$, --- непрерывные, достаточно гладкие функции,


удовлетворяющие условиям 


\begin{gather}\label{2.3}


\vf^1(x,t)+\vf^2(x,t)=\vf(x,t),\quad (x,t)\in S,\quad


\vf^2(x,t)\equiv 0,\quad r(x,\Gamma)\ge \de,\;(x,t)\in S,\notag\\


%


\left|\dxx \vf^1(x,t)\right|\le M,\quad (x,t)\in S_0,\quad k\le 6;\quad


\left|\dxx \vf^2(x,t)\right|\le M\de^{-k},\quad (x,t)\in S_0,\tag{2.3}


\end{gather}


где $r(x,\Gamma)$ --- расстояние от точки $x$ до множества $\Gamma$.


Запись $L_{(j.k)}$ (или $f_{(j.k)}(x,t)$) означает, что эти операторы 


(функции) впервые введены в формуле $(j.k)$.





Для решения краевой задачи (1.1) и его компонент из представления (2.1)


справедлива оценка


\begin{align}\label{2.4}


&\left|\dxt u(x,t)\right| \le M \de^{-k-k_0}\left[ 1+


\eps^{2k+2k_0}\de^{-2k-2k_0}\right], \quad (x,t)\in \ovl{G},


\tag{2.4а}\\


%


\intertext{а также оценки}


%


&\left|\dxt U(x,t)\right| \le M,\tag{2.4б}\\


%


&\left|\dxt V(x,t)\right| \le M \de^{-k-k_0}\left[ 1+


\eps^{2k+2k_0}\de^{-2k-2k_0}\right], \quad (x,t)\in \ovl{G}, \quad


k+2k_0\le 4.\notag


\end{align}





Для вычисления сингулярной части решения $V(x,t)$ будет удобно использовать 


запись краевой задачи в новых переменных. От переменных $x,$ $t$ перейдем


к переменным $\eta,$ $t$ 


$$


\eta=\eta(x,t)=x-Q(t), \quad (x,t)\in\ovl{G}.


$$


Здесь соотношение 


$$


x=Q(t),\quad t\in [0,T],\quad Q(0)=0,


$$


задает характеристику вырожденного уравнения (уравнения (1.1) при 


$\eps=0$), являющуюся решением дифференциального уравнения


$$


\frac{dx}{b(x,t)} = - \frac{dt}{p(x,t)}


$$


и проходящую через точку $x=0,$ $t=0$. Нетрудно проверить, что 


при переходе к новой (характеристической) переменной $\eta$ 


дифференциальный оператор $L_{(1.1)}$ преобразуется в оператор, 


не содержащий конвективного члена \(b(x,t)(\pa / \pa x)\).





Через $\wh{V}_0(\eta,t)$, $(\eta,t)\in\ovl{\wh{G}}$, обозначим решение


вспомогательной задачи 


\begin{align} \label{2.5}


L_{(2.5)}\widehat{V}_0(\eta,t) &\equiv \left\{\eps^2\widehat{a}(0,t)


\frac{\pa^2}{\pa\eta^2}-\widehat{c}(0,t)-


\widehat{p}(0, t)\frac{\pa}{\pa t}


\right\}\widehat{V}_0(\eta,t)=0,\quad 


(\eta,t)\in\widehat{G},\tag{2.5}\\


\widehat{V}_0(\eta,t)&=\widehat{\vf}^2(\eta,t),\quad 


(\eta,t)\in \widehat{S}.\notag


\end{align}


Здесь $\widehat{G}^0=\{(\eta,t):\eta=\eta(x,t),\ (x,t)\in G^0\}$ 


--- образ множества $G^0$, $G^0\subset \ovl{G}$,


$\wh{v}(\eta,t)=v(x(\eta,t),t)$,


где $v(x,t)$ есть одна из функций $a(x,t),$ $c(x,t),$ $p(x,t),$


$\vf^2(x,t)$. Взяв коэффициенты и правую часть дифференциального 


уравнения (2.5) вдоль характеристики $x=Q(t)$, мы еще более упростили


определение сингулярной части решения.





Заметим, что регулярная (гладкая) компонента из представления 


(2.1) является решением краевой задачи с достаточно гладкими данными


(см. задачу (2.2а)). При малых значениях параметров $\eps$ и 


$\delta$ сингулярная часть решения задачи (см. задачу (2.2б)) 


достаточно мала вне малой окрестности характеристики, проходящей через


середину интервала, на котором начальная функция претерпевает возмущение.


Это позволяет сингулярное решение приблизить решением более простой 


задачи --- краевой задачи (2.5) для дифференциального уравнения с


коэффициентами, ``замороженными'' в окрестности упомянутой выше 


характеристики. Уравнение (2.5) уже не содержит конвективных членов. 


Для таких уравнений построение подходящих сеточных аппроксимаций


уже рассматривалось (напр., \cite{Sh2}).





При построении специальных разностных схем (см. \S~4) для вычисления


сингулярной компоненты $V(x,t)$ используются специальные сетки, сгущающиеся


в той подобласти, где сингулярное решение существенно. Это приводит


к упрощению численного метода.





\section{Классическая разностная схема}





Рассмотрим классическую разностную схему для решения задачи (1.1). В качестве 


классической выберем монотонную разностную схему А.А.Самарского 


(\cite{Sa} гл.\,VII, \S\,1, с.\,401). 


На множестве $\ovl{G}$ введем сетку 


\begin{equation}\label{3.1}


\ov{G}_h=\ov{D}_h\times\ovo_0=\ov{\om}\times\ovo_0,


\end{equation}


где $\ovo,$ $\ovo_0$ --- сетки на интервалах $[-d,d]$ и $[0,T]$, 


вообще говоря, неравномерные.


Полагаем $h^i=x^{i+1}-x^i,$ $x^i,\;x^{i+1}\in\ovo,$


$h^k_t=t^{k+1}-t^k,$ $t^k,\;t^{k+1}\in\ovo_0$; 


$h=\max_{i}h^i,$ $h_t=\max_{k}h^k_t$.


Через $N+1$ и $N_0+1$ обозначим число узлов сеток $\ovo$ и $\ovo_0$ 


соответственно; полагаем $h\le MN^{-1},$ $h_t\le MN_0^{-1}$.


Будем рассматривать также сетки 


\begin{equation}\label{3.2}


\ov{G}_h=\ov{G}_{h(\ref{3.1})},


\end{equation}


у которых сетка по времени $\ovo_0$ равномерная.


 


На сетке $\ov{G}_{h(\ref{3.1})}$ аппроксимируем


краевую задачу (\ref{1.1}) разностной схемой 


\begin{align} \label{3.3}


\begin{split}


\Lambda_{(\ref{3.3})}z(x,t)&=f(x,t),\quad (x,t)\in G_h,\\


z(x,t)&=\vf(x,t),\quad (x,t)\in S_h.


\end{split}


\end{align}


Здесь $\ov{G}_h=\ov{G}_{h(\ref{3.1})}$,


$$


\Lambda_{(\ref{3.3})}z(x,t)\equiv\left\{\,\eps^2 a(x,t)


\de_{\ov{x}\,\widehat{x}}+


b^+(x,t)\de_{x}+b^-(x,t)\de_{\ov{x}} - 


c(x,t)-p(x,t)\de_{\ov{t}}\,\right\}z(x,t),


$$


$b^+=2^{-1}(b+|b|)\ge 0$, $b^-=2^{-1}(b-|b|)\le 0$,


$\de_{x\,} z(x,t)$, $\de_{\ov{x}}\,z(x,t)$ и


$\de_{\ov{x}\,\widehat{x}}\,\,z(x,t)$ --- первые (прямая и обратная) и


вторая разностные производные на неравномерных сетках, 


$\de_{\ov{t}}\,z(x,t)$ --- обратная разностная производная по времени,


например,


$$


\de_{\ov{x}\,\widehat{x}}\,\,z(x,t)=2\left(h^{i-1}+h^i\right)^{-1}


(\de_x\,z(x,t)-\de_{\ov{x}}\,z(x,t)), \quad x=x^i.


$$





В силу принципа максимума решение разностной схемы (3.3), (3.1)


ограничено $(\eps,\delta)$-равномерно:


$|z(x,t)| \le M$, $(x,t)\in\ovl{G}_h$.


Исследование сходимости схемы 


проводится подобно исследованию схем в \cite{Sh1}, \cite{Sh5}. 





Принимая во внимание априорные оценки вида (2.4), с помощью 


принципа максимума устанавливаем сходимость разностной схемы (3.3), (3.1) 


при фиксированном $\delta$


\begin{align}\label{3.4}


|\,u(x,t)-z(x,t)\,|&\le M[\eps^8\de^{-9}N^{-1}+\eps^4\de^{-6}


N_0^{-1}],\quad \eps\ge\de;\\


%


|u(x,t)-z(x,t)|&\le M[\delta^{-1}N^{-1}+\delta^{-2} N_0^{-1}],\quad 


\eps\le \de,\quad (x,t)\in\ov{G}_h.\notag


\end{align}





\begin{thmm}\label{th1}


Пусть решение краевой задачи {\em (\ref{1.1})} достаточно гладко при 


фиксированных значениях возмущающих параметров. 


Тогда решение разностной схемы {\em (\ref{3.3}), (\ref{3.1})}


сходится $\eps$-равномерно при фиксированных значениях параметра


$\delta;$ для него справедливы оценки {\em (\ref{3.4}).}


\end{thmm} 





\section{Специальная разностная схема для задачи \protect(\ref{1.1})}


 


{\bf 4.1.} Для построения специальной разностной схемы, сходящейся 


($\eps,\delta$)-равномерно, будем использовать метод декомпозиции


решения в окрестности локального возмущения начального условия \cite{Sh2}. 


Решение задачи аппроксимируем функцией


$$


u^0(x,t)=U(x,t)+V_0(x,t),\quad (x,t)\in\ovl{G}.


$$


Здесь $U(x,t)=U_{(\ref{2.1})}(x,t)$ --- решение задачи (2.2a),


$V_0(x,t)=\wh{V}_{0(\ref{2.5})x}(x,t)=\wh{V}_{0(2.5)}(\eta(x,t),t)$,


$(x,t)\in\ovl{G}$, --- решение вспомогательной задачи (2.5), записанное


в исходных переменных $x,\ t$. Функция $V_0(x,t)$ является главным членом


асимптотического разложения сингулярной части решения краевой задачи.


При малых значениях параметра $\delta$ функции


$V(x,t)$ и $V_0(x,t)$ близки. Для функции $V_0(x,t)$, $(x,t)\in \ovl{G}$,


справедливы оценки, подобные приведенным в \cite{Sh2}.





Краевую задачу (2.2a) (для гладкой части $U(x,t)$) 


аппроксируем классической разностной схемой


\begin{align}\label{4.1}


\begin{split}


\Lambda_{(\ref{3.3})}z(x,t)&=f(x,t),\quad (x,t)\in G_h,\\


z(x,t)&=\vf^1(x,t),\quad (x,t)\in S_h,


\end{split}


\end{align}


где $\ov{G}_h=\ov{G}_{h(\ref{3.2})}$.





Построим разностную схему для задачи (2.5).


На множестве $\ovl{\wh{G}}$ с криволинейной границей $\wh{S}$ на плоскости


($\eta,t$) введем сетку 


\begin{equation}\label{4.2}


\ovl{\wh{G}}_h=\widehat{G}_h\cup\widehat{S}_h,\quad 


\widehat{G}_h=\widehat{G}\cap


\left\{\om\times\ovo_0\right\},


\end{equation}


где $\ovo_0=\ovo_{0(\ref{3.1})},$ $\om$ --- сетка на оси $\eta$, 


вообще говоря, неравномерная. Через $N+1$ обозначим число узлов сетки


$\om$ на минимальном отрезке оси $\eta$, на который проектируется


множество $\ovl{\wh{G}}$. Множество $\widehat{S}_h$ состоит из точек 


пересечения границы $\widehat{S}$ с прямыми, проходящими через узлы 


множества $\widehat{G}_h$ параллельно осям.





На сетке $\ovl{\wh{G}}_h$ задаче (2.5) сопоставим разностную схему


\begin{align}\label{4.3}


\begin{split}


\Lambda_{(\ref{4.3})}\widehat{z}(\eta,t)&= 0,\quad


(\eta,t)\in \widehat{G}_h,\\


\widehat{z}(\eta,t)&=\widehat{\vf}^2(\eta,t),\quad


(\eta,t)\in \widehat{S}_h,


\end{split}


\end{align}


где 


$$


\Lambda_{(\ref{4.3})}\widehat{z}(\eta,t)\, \equiv 


\{\eps^2 \widehat{a}(0,t)\de_{\ov{\eta}\widehat{\eta}}-


\widehat{c}(0,t)-\widehat{p}(0,t)\de_{\ov{t}}\}


\widehat{z}(\eta,t).


$$





Приближенное решение задачи (\ref{1.1}) определим соотношением


\begin{equation}\label{4.4}


z(x,t)=\ov{z}_{(\ref{4.1},\ref{3.2})}(x,t)+


\widehat{z}_{(\ref{4.3},\ref{4.2})\,x}(x,t),\quad


(x,t)\in \ov{\widehat{G}}_{h\,x}.


\end{equation}


Здесь $\ovl{\wh{G}}_{h\,x}=\{(x,t):x=x(\eta,t),\; (\eta,t)\in 


\ovl{\wh{G}}_h\}$ --- образ множества $\ovl{\wh{G}}_{h}$ при отображении 


$x=x(\eta,t)\equiv \eta+Q(t),$ $(\eta,t)\in \ov{\widehat{G}}$;


$\widehat{z}_{(\ref{4.3},\ref{4.2})x}(x,t)=


\widehat{z}_{(\ref{4.3},\ref{4.2})}(\eta(x,t),t)$ --- образ решения 


задачи (\ref{4.3}), (\ref{4.2}). Через 


$\ov{z}_{(\ref{4.1},\ref{3.2})}(x,t)$ обозначаем билинейный интерполянт


решения задачи (\ref{4.1}), (\ref{3.2}).





Для решения задачи (2.5) используем сетки, сгущающиеся в 


окрестности множества $\widehat{S}^*$, где


$\wh{S}^*=\{(x,t):x=Q(t),\;t\in(0,T]\}$.





{\bf 4.2.} Пусть выполняется условие


$$


\eps\le \de,\quad \eps\in(0,1],\quad \de\in(0,d].


$$


Сначала зададим на множестве $\ovl{G}$ специальную сетку, сгущающуюся 


в окрестности точки $x=0,$


\begin{equation}\label{4.5}


\ov{G}_h=\ov{D}_h^с\times\ovo_0=\ovo^с\times\ovo_0,


\tag{4.5а}


\end{equation}


где $\ovo_0=\ovo_{0(\ref{3.2})},$ $\ovo^с=\ovo^с(\si)$ ---


кусочно-равномерная сетка на $[-d,d]$, $\si$ --- параметр, зависящий от 


$\de$ и $N$.


Сетка $\ovo^с$ строится равномерной на каждом из отрезков $[-\si,\si]$ 


и $[-d,-\si],$ $[\si,d]$ с шагом $h^{(1)}=4\si N^{-1}$ и 


$h^{(2)}=4(d-\si)N^{-1}$ соответственно.


Величину $\si$ выбираем удовлетворяющей условию 


\begin{equation}


\si=\si_{(\ref{4.5})}(\de,N)=\min\left[2^{-1}d,\;m^{-1}\de\ln N\right],


\tag{4.5б}


\end{equation}


где $m>0$ --- произвольное число.





На множестве $\ovl{\wh{G}}$ введем специальную сетку 


\begin{equation} \label{4.6}


\ovl{\wh{G}}_h=\ovl{\wh{G}}_{h(\ref{4.2})},\tag{4.6}


\end{equation}


где $\om=\om^с(\si)$ --- кусочно-равномерная сетка, подобная сетке


$\ov{\om}^{\,с}_{(\ref{4.5})}$, сгущающаяся в окрестности точки 


$\eta=0$. Считаем распределение узлов сетки $\om^с_{(4.6)}(\si)$ 


на отрезке $[-\si,\si]$ и вне его равномерным; шаги $h^{(1)}$ и 


$h^{(2)}$ соответственно на отрезке $[-\si,\si]$ и вне его определяем


соотношениями: $h^{(1)}=4\sigma N^{-1}$, $h^{(2)}=2(d_{\eta}-2\sigma) N^{-1}$,


где $\si=\si_{(\ref{4.5})}(\de,N)$, $d_{\eta}$ --- длина отрезка на оси 


$\eta$, на который проектируется множество $\ovl{\wh{G}}$.





Для решения разностной схемы (\ref{4.1}), (\ref{3.2}), 


(\ref{4.3}), (\ref{4.6}) получаем оценку 


$$


\left|u(x,t)-z(x,t)\right|\le M\left[N^{-1}\ln N+ N_0^{-1}+\de\right],\quad


(x,t)\in\ovl{\wh{G}}_{h\,x},\quad


\ovl{\wh{G}}_h=\ovl{\wh{G}}_{h(4.6)},\quad \eps\le\de.


$$





В том случае, если выполняется условие 


$\eps\ge \de$, $\eps\in(0,1]$, $\de\in(0,d]$,


строим на $\ovl{\wh{G}}$ сетку, сгущающуюся в окрестности множества 


$\widehat{S}^*$, а также при малых значениях $t$


\begin{equation}\label{4.7}


\ovl{\wh{G}}_h=\ov{G}_{h(\ref{3.1})} \ \{\text{ при } \ 


(x,t)\equiv (\eta,t)\ \},\tag{4.7}


\end{equation}


где $\om=\om^c_{(4.6)}(\si)$ при 


$\si=\si_{(\ref{4.5})}(\de,N)$;


$\ovo_0=\ovo^c_0(\si_0)$ --- кусочно-равномерная сетка на отрезке $[0,T]$, 


$\si_0$ --- параметр, зависящий от $\eps$, $\de$, $N$ и $N_0$. Определим 


величину 


$\si_0=\min\,[\,2^{-1}T,\eps^{-2}\de^2 N_1^{2/3}, 


\eps^{-2}\de^2 N_0^{2/5}]$. Шаги сетки $\ovo^c_0$ на интервалах 


$[0,\si_0]$ и $[\si_0,T]$ постоянны и равны $h_0^{(1)}=2\si_0 N_0^{-1}$ и


$h_0^{(2)}=2(T-\si_0)N_0^{-1}$ соответственно. 


При $\eps=\de$ сетка (4.7) совпадает с сеткой (4.6). 





Для решения разностной схемы 


(\ref{4.1}), (\ref{3.2}), (\ref{4.3}), (4.7)


получается оценка 


$$


\left|u(x,t)-z(x,t)\right|\le M\left[N^{-1/3}+ N_0^{-1/5}


+\de^{1-\nu}\right],\quad


(x,t)\in\ovl{\wh{G}}_{h\,x},\quad \ovl{\wh{G}}_h=


\ovl{\wh{G}}_{h(4.7)},\quad \eps\ge\de,


$$


где $0<\nu<1$ --- произвольное сколь угодно малое число.





Таким образом, при условии $\eps\le\delta$ (или $\eps\ge\delta$)


разностная схема (4.3) на специальной сетке (4.6) (сетке (4.7))


сходится к решению краевой задачи (\ref{1.1}) 


при $N,\, N_0\to\infty$ для $\delta=o(1)$.





{\bf 4.3.} Построим $(\eps,\de)$-равномерную аппроксимацию 


решения краевой задачи (\ref{1.1}). 


Определим функцию $z_{(4.8)}(x,t),$


$(x,t)\in \ov{G}_h$, --- приближенное решение краевой задачи, полагая 


\begin{equation}\label{4.8}


z(x,t)=z_{(\ref{3.3},\ref{3.2})}(x,t),\quad 


(x,t)\in \ov{G}_{h(\ref{3.2})},\tag{4.8}


\end{equation}


$$


\text{при } \eps\le \de,\; \de\ge \Psi_1(N,N_0) \text{ либо } 


\eps>\de, \; \de \ge \Psi_2(N,N_0);


$$


$$


z(x,t)=z_{(\ref{4.1},\ref{3.2},\ref{4.3},\ref{4.6})}(x,t),


\quad (x,t)\in \ov{\wh{G}}_{h(4.6)x},


$$


$$


\text{при } \eps\le \de,\ \ \de<\Psi_1(N,N_0);


$$


$$


z(x,t)=z_{(\ref{4.1},\ref{3.2},\ref{4.3},(4.7)}(x,t),\quad 


(x,t)\in \ov{\wh{G}}_{h(4.7)x},


$$


$$


\text{при } \eps>\de,\ \ \de<\Psi_2(N,N_0),


$$


где 


$$


\Psi_1(N,N_0)=[N^{-1}+N_0^{-1}]^{1/3}, \quad


\Psi_2(N,N_0)=[N^{-1}+N_0^{-1}]^{1/10}.


$$





Решение специальной разностной схемы сходится 


$(\eps,\de)$-равномерно; для сеточного решения выполняются оценки 


\begin{align}\label{4.9}


|u(x,t)-z_{(4.8)}(x,t)|&\le M \left(N^{-1}


+N_0^{-1}\right)^{1/3},\quad


(x,t)\in \ov{G}_h,\quad \eps\le\de;\tag{4.9}\\


%


|u(x,t)-z_{(4.8)}(x,t)|&\le M 


\left(N^{-1}+N_0^{-1}\right)^{(1-\nu)/10},\quad


(x,t)\in \ov{G}_h,\quad \eps>\de.\notag


\end{align}





\begin{thmm} \label{th2}


Пусть решение краевой задачи {\em (\ref{1.1})} достаточно гладко при 


фиксированных значениях параметров. Тогда приближенное решение 


$z_{(4.8)}(x,t)$, $(x,t)\in \ov{G}_h$, 


сходится к решению краевой задачи $(\eps,\de)$-равномерно.


Для решения специальной разностной схемы справедливы оценки {\em (4.9).}


\end{thmm}





\section{Численный пример}





Эффективность предложенного в \S\,4 алгоритма иллюстрирует


модельная задача


\begin{equation}\label{5.1}


\eps^2\frac{\partial^2 u}{\partial x^2}+\frac{\partial u}{\partial x}


=\frac{\partial u}{\partial t},\quad 0<x<1,\quad 0<t\le 1,


\tag{5.1а}


\end{equation}


%


\begin{equation}


u(0,t)=u(1,t)=0,\quad u(x,0)=\vf_0(x),


\tag{5.1б}


\end{equation}


где 


$$


\vf_0(x)=


\begin{cases}


0, & |x|\ge \delta;\\


(1-x/\delta)^5\,(1+x/\delta)^5, & |x|\le \delta.


\end{cases}


$$


Отметим, что данные краевой задачи (5.1) подобраны так, чтобы


пограничные слои не возникали. 





По схемам, предложенным в \S\;\S\;3,\;4, была проведена серия расчетов с 


различными значениями параметров $\eps$ и $\de$ ($\eps, \delta=4^{1-k},$ 


$k=1,\dotsc,5$) и различным числом узлов $N$ разностной сетки 


($N=4^k,$ $k=1,\dotsc ,4$) при условии $N=N_0$. Для оценки эффективности разностных 


схем сравнивали величину погрешности приближенного решения,


определяемую по формуле


$$


E(N=N_0,\eps,\delta)=\max_{\ovl{G}_h}\, |u^{\star}(x,t)-z(x,t)|.


$$


Здесь $u^{\star}(x,t)$ есть билинейная интерполяция по $x$ и $t$ функции


$z^{\star\,512}(x,t)$, являющейся решением специальной разностной схемы на сетке с


числом узлов $N=N_0=512$.





Проведенные расчеты показали, что классическая 


разностная схема на равномерных сетках дает значения решения задачи (5.1) 


с большой ошибкой при малых значениях $\delta$. Анализируя поведение 


погрешности, мы пришли к выводу, что дальнейшее уменьшение параметров 


$\eps$ и $\de$ ведет к полной потере точности. Тогда как при использовании


специальной разностной схемы ошибка стабилизируется при уменьшении 


$\eps$ и $\de$. При этом с ростом числа узлов расчетной


сетки точность решения растет.


Таким образом, предложенный сеточный метод позволяет 


получить решение с достаточной точностью, практически не зависящей от 


величины возмущающих параметров. 





\section{Замечания и обобщения}





{\bf 6.1.} В том случае, когда решение краевой задачи (1.1) содержит


пограничные слои, при построении специальных разностных схем используем 


сетки $\ovl{G}_h$, сгущающиеся также в окрестности множества


$S_1$ --- боковой границы области $G$. Сетку $\ovl{\omega}=\ovl{\omega}^{\,c}$,


используемую при конструировании $\ovl{G}_h$, строим кусочно


равномерной; закон сгущения узлов в окрестности концов отрезка


$[-d,d]$ выбираем согласно \cite{Sh1}, \cite{Sh4}. Для схем, построенных таким


образом, сохраняются утверждения теоремы о сходимости равномерно


относительно параметров $\eps$ и $\delta$.





{\bf 6.2.} При обосновании сходимости специальных схем предполагались


выполненными априорные оценки решения и его производных, а также оценки 


регулярной и сингулярной частей решения, установленные в случае условий


(1.2) и (2.3). При нарушении этих условий начальные


данные приближаем в равномерной метрике подходящими гладкими функциями,


удовлетворяющими указанным условиям. Таким образом, удается показать


равномерную относительно параметров $\eps$ и $\delta$ сходимость


построенных специальных разностных схем, однако порядок сходимости в


этом случае снижается.





 


Авторы выражают свою признательность П.А.Фарреллу (Кент, США) и


П.В.Хемкеру (Амстердам, Нидерланды) за полезные обсуждения и замечания,


стимулирующие эти исследования.








\begin{thebibliography}{99}





\bibitem{Ba}


Бахвалов Н.С. {\em К оптимизации методов решения краевых задач


при наличии пограничного слоя} // Журн. вычисл. матем. и матем. физ. --


1969. -- Т.\,9. -- \No\,4. -- С.\,841--859. 





\bibitem{Il}


Ильин A.M. {\em Разностная схема для дифференциального уравнения с малым 


параметром при старшей производной} // Матем. заметки. -- 1969. --


T.\,6. -- Вып.\,2. -- C.\,237--248. 





\bibitem{Du}


Дулан Э., Миллер Дж., Шилдерс У. {\em Равномерные численные методы 


решения задач с пограничным слоем.} -- M.: Мир, 1983. -- 199 c.





\bibitem{Sh1}


Шишкин Г.И. {\em Сеточные аппроксимации сингулярно возмущенных 


эллиптических и параболических уравнений.} -- Екатеринбург: УрО РАН, 1992. 


-- 233 с.





\bibitem{Mar}


Марчук Г.И., Шайдуров В.В. {\em Повышение точности решений 


разностных схем.} -- М.: Наука, 1979. -- 319 с.





\bibitem{Sh2} 


Шишкин Г.И. {\em Сеточная аппроксимация параболических уравнений 


с сингулярными начальными условиями} // Журн. вычисл. матем. и 


матем. физ. -- 1996. -- Т.\,36. -- \No\,3. -- С.\,73--92.





\bibitem{Ti}


Тихонов А.Н., Самарский А.А. {\em Уравнения математической физики.} --


3-е изд. -- М.: Наука, 1966. -- 724 с.





\bibitem{Sh3}


Шишкин Г.И. {\em Проблема аппроксимации диффузионного потока при


численном моделировании процесса переноса примесей} // Матем.


моделир. -- 1995. -- Т.\,7. -- \No\,7. -- С.\,61--80.





\bibitem{Sh4}


Шишкин Г.И. {\em Разностная аппроксимация сингулярно возмущенной 


краевой задачи для квазилинейных эллиптических уравнений, 


вырождающихся в уравнение первого порядка} // Журн. вычисл. матем. и 


матем. физ. -- 1992. -- Т.\,32. -- \No\,4. -- С.\,550--566.





\bibitem{Sh5}


Shishkin G.I. {\em Grid approximation of a singularly perturbed 


boundary value problem for a quasi-linear parabolic equations in the case 


of complete degeneracy in spatial variables} // Sov. J. Numer. Anal. and


Math. Model. -- 1991. -- V.\,6. -- \No\,3. -- C.\,243--261.





\bibitem{Sa}


Самарский А.А. {\em Теория разностных схем.} -- 3-е изд. -- 


М.: Наука, 1989. -- 616 с.





\end{thebibliography}





\vskip 2cm





\post{\it Институт математики и механики}{\it Поступила}


\post{\it Уральского отделения}{01.10.1996}


\post{\it Российской Академии Наук}{}





\end{document}





